% !TeX root = ../main.tex
\subsection{Model Choice and Mathematical Rationale}
\label{sec:model-choice}

The model must learn from sparse, irregularly spaced results, update both opponents, distinguish draws, and represent uncertainty before a result is known. It must also separate team ability from the advantage of playing at home. These requirements determine which parts of earlier work we use. Elo-derived football predictors provide a useful separation between ratings and outcome regression~\cite{hvattum2010elo}, but a fixed-step point rating does not itself describe uncertainty about ability. We retain a fixed-step internal baseline for comparison, rather than interpreting the numerical difference between two rating points as an observed physical difference in football quality. Our latent difference is a coordinate in a probability model; its unit is fixed by performance-noise variance, and its effect on probability is nonlinear.

Dynamic paired comparisons~\cite{glickman1999dynamic} supply the central distinction between latent ability and uncertainty about that ability. A Gaussian evolution model is a tractable starting point: inactivity enlarges uncertainty while leaving the expected strength unchanged. However, a single smooth diffusion may adapt poorly to persistent abrupt changes. We therefore develop both the smooth special case and a model with an additional persistent jump component. The jump model is selected for the admitted development requirements, not because mathematical complexity guarantees better prediction. A simpler model can remain preferable in a stationary simulated setting.

TrueSkill~\cite{herbrich2006trueskill} offers related ideas of uncertain skill, noisy performance, and a draw region. These ideas suit a result archive containing home wins, draws, and away wins. Its game-rating application does not by itself establish validity for national teams. Our dated evolution, venue term, robust likelihood, and chronological checks are additional project decisions. Expectation propagation~\cite{minka2001ep} supplies an approximate way to communicate outcome information through correlated latent states; it is a computational method, not a guarantee that the entire posterior is Gaussian or exact.

A goal-count model, illustrated by Dixon and Coles~\cite{dixon1997scores}, could use more of the recorded score information. We do not adopt it for this analysis: the principal outcome is a home win, and scores are not uniformly restricted to regulation time. Adding a score model would require its own observation assumptions and evaluation rather than treating goals and the resulting win as independent evidence. FIFA rankings and published Elo ratings are not tested external benchmarks here. Neither the present comparisons nor the cited papers establish that our implementation outperforms those systems.

Two other method choices concern different problems. Direct standardization answers how observed rates differ under a common match composition; it does not learn team strength. Proper scoring rules~\cite{gneiting2007scoring} evaluate whether probabilities are useful for subsequent outcomes, while simulation-based uncertainty analysis examines stated comparison procedures under specified generating mechanisms. No one of these tools substitutes for the others. In particular, model-conditional interval coverage cannot establish real-world predictive superiority.
